\section{Europe’s Spirit in the North}

Gornahoor has deliberately underscored the ``Russian connection''\footnote{\url{https://gornahoor.net/?p=334}} for traditionalism in the previous century (and before). As if to emphasize such a connection explicity ``Romane'', the reception of Evola via Samizdat circles in the Soviet era demonstrates the resonance virtually all of his ideas\footnote{\url{http://www.amerika.org/tag/alexandre-douguine/}} had with the ideas of those in Russia who were not linked with the monarchy by tradition, but had ``come into'' it as part of the new Right.

This is all the more remarkable when one considers the circles that Evola kept. It also occurred in correlation with influence from Gurdjieff\footnote{\url{http://traditionalistblog.blogspot.com/2007/03/russian-traditionalism-started-with.html}}, because these certain individuals who basically lived in the Lenin Library were a ``fourth way'' group:

\begin{quotex}
Stepanov, Golovin, Haydar Jamal, Yuri Mamleyev and Rovner were among those who for a few years ``almost lived'' in the Lenin Library, not only reading but also eating and--above all--talking there. They formed a small and very close Gurdjieff group--or, Rovner insists, a ``company,'' as its members were ``too individualist'' to form a group.
\end{quotex}


It was an offshoot of the Gurdjieff group which had assimilated Guenon's corpus \& eventually stumbled over Evola.

Of course, to the modern West, these persons are all lunatics\footnote{\url{http://ku-eichstaett.academia.edu/AndreasUmland/Papers/197791/Aleksandr\_Dugins\_Transformation\_from\_a\_Lunatic\_Fringe\_Figure\_into\_a\_Mainstream\_Political\_Publicist\_19801998\_A\_Case\_Study\_in\_the\_Rise\_of\_Late\_and\_Post-Soviet\_Russian\_Fascism}}, and very prestigious individuals receive lots and lots of money to research them in order to expose their ``rise to power''.

A debate between Aleksandr Dugin \& a Brazilian interlocutor naturalized to the USA is quite revealing. In broken English\footnote{\url{http://www.theinteramerican.org/blogs/olavo-de-carvalho/275-olavo-de-carvalho-debates-alexandr-dugin-iii.html}}, Dugin attempts to convey his dismay and ultimate perspective on the liberal vs. conservative struggle for the soul of ``America'':

\begin{quotationx}
It is obvious that modern Eastern societies have many negative aspects. But they are mostly the result of modernization, westernization and the perversion of the ancient traditions. In my youth (early 80-s) I was anticommunist in the Guenonian/Evolian sense. But after having known modern Western Civilization and especially after the end of Communism \emph{I have changed my mind} and revised this traditionalism discovering the other side of the socialist society, which is the parody on the true Tradition, but nevertheless is \emph{much better than absolute absence} of the Tradition in Modern and Post-Modern Western world. So, I love the East in general and blame the West. The West now expands itself on the planet. So the globalization is Westernization and Americanization. Thereforee, I invite all the rest to join the camp and fight Globalism, Modernity/Hypermodernity, Imperialism Yankee, liberalism, free market religion and unipolar world. These phenomena are the ultimate point of the Western path to the abyss, the final station of the evil and the almost transparent image of the antichrist/ad-dadjal/erev rav. So the West is the center of kali-yuga, its motor, its heart. \emph{Mr. Carvalho blames the East and loves the West}. But here begins some asymmetry.

\textbf{I love the East as a whole including its dark sides. The love is the strong, very strong feeling. You don't love only good and pure sides of the beloved one, you love him wholly. Only such love is real one. Mr. Carvalho loves the West but not all the West, only its part. The other part he rejects.}

To explain his attitude in front of the East he makes appeal to the conspiracy theory. Scientifically it is inadmissible and discredits immediately Mr. Carvalho thesis but in this debate I don't think that scientific correctness is that does mean much. I don't try to please or convince somebody. I am interested only in the truth (\emph{vincit omnia veritas}). If Mr. Carvalho prefers to make use of the conspiracy theory let him do it. The Western Christianity stressed the individual as the center of the religion and made the salvation the strictly individual affair.

\textbf{The Protestantism led this tendency to the logical end. Denying more and more the holistic ontology of the organic society the Western Christianity arrived with the Modernity to self-denial (deism, atheism, materialism, economism). French sociologist Luis Dumont in his excellent books \emph{Essai sur l'Individualism} and \emph{Homo Aequalis} shows that the methodological individualism is the result of the oblivion and direct purge by the Western scholastic of the early and original Greco-Roman theological tradition conserved intact in the Byzance and Eastern Church as whole... My opinion: American paleoconservatives, traditional American right are doomed. Their discourse is incoherent, weak and too idiosyncratic}.
\end{quotationx}


Dugin's politics may or may not be ``proto-fascistic'', \& Carvalho may or may not have a good response, but Dugin certainly grasps the concept that love emanates from something other than a knee-jerk \& shallow hatred of ``the other'' who disagrees over man's right to strip himself metaphysically.

The occurrence of so many threads in Russia is intriguing to those interested in ``coincidences'' and the course of Empires. If the star of the West has moved consistently Westward (and no less a secular scholar than Ortega y Gasset took such a theory seriously), it has crossed back into Asia from the Pacific (Gasset, by the way, thought that the disappearance of oak trees in Europe heralded something real and ominous). Russia may be the reconciliation point between the East \& the West; Russia rather than China.

\emph{Homo Atlanticus}, once more, is on the wane, even as he approaches his zenith. Holy Mother Russia unaccountably remains.

PA Sorokin, an emigre to America who founded Harvard's Sociology Department, used his cyclical view of history to locate the West:

\begin{quotex}
\emph{Ideational truth} is the truth revealed by the grace of God, though his mouthpieces (the prophets, mystics, and founders of religions), disclosed in a supersensory way through mystic experience, direct revelation, divine intuition, and inspiration. Such a truth may be called the truth of faith. It is regarded as infallible, yielding adequate knowledge about the true-reality values. \emph{Sensate truth} is the truth of the senses, obtained through our organs of sense perception.
\end{quotex}


He goes on to argue that \emph{idealistic truth} is a synthesis of what is known by Faith \& what is known by empirical experience. One could sharpen this insight\footnote{\url{http://www.gornahoor.net/?p=1189}} by pointing to the union of sense \& faith in Evola's heightened metaphysical positivism, or magical realism. Such an \emph{apotheosis} in the individual obviates the need to maintain the delicate balance between ideational truth and sensate experience; it (as it were) sets the individual outside of the cycles of time \& space.

Russia, set as it is as a potential adjunct or anteroom of Europe, with access to all of her high culture, but straddling the waste places of Asia and the deserts of the East, above all living in the cold and the north, is admirably suited to become the arbiter of a Third Romanity. The continual wrestling match which this country has engaged Modernity \& Enlightenment since inception augurs well for its future stars.

Mouravieff's explication of Peter the Great's providential role\footnote{\url{http://www.bibliotecapleyades.net/archivos\_pdf/mouravieff3.pdf}} in raising up Russia p.19 (as well as the importance of the Hellenic Orthodox world to universal history, along with speculations about China's future) are also important to keep in mind.


\flrightit{Posted on 2011-06-18 by Logres}

\begin{center}* * *\end{center}

\begin{footnotesize}
\begin{sffamily}
\texttt{Jason on 2011-06-30 at 20:12 said:}

I'm surprised that no one on this site has yet mentioned Dmitri Merejkowski, his books are on the very same path that Gornahoor is treading. Basing himself on the Stromata of Clement, Merejkowski created a philsophical and aesthetic system that saw Paganism as Christianity before Christ and the pagan Gods as shadows of the coming Lord... ... 

\hfill

\end{sffamily}
\end{footnotesize}

